\documentclass[aspectratio=169]{beamer}

\usepackage{amsmath}
\usepackage{amssymb}
\usepackage{booktabs}
\usepackage{array}

\usetheme{Madrid}
\setbeamertemplate{navigation symbols}{}

\title{Mission-to-Execution Runtime Validation}
\subtitle{Runtime Modeling for Multi-Robot Systems}
\author{Huan Zhang}
\date{\today}

\begin{document}

\begin{frame}
\titlepage
\end{frame}

\begin{frame}{Research Problem}
The offline SMT model produces a valid planned execution:
\[
x^{plan} \models \Phi_{\mathrm{global}}
\]
from which the execution graph $G_E$ is constructed.

However, during execution,
\[
x^{obs}(t) \neq x^{plan}
\]
because of task delays, robot failures, travel-time variations, or resource contention.

\begin{center}
\boxed{
\begin{minipage}{0.82\textwidth}
\centering
Does the observed execution still satisfy the mission guarantees?
\end{minipage}
}
\end{center}
\end{frame}

\begin{frame}{Why Runtime Validation?}
The SMT solution guarantees that the planned execution satisfies the mission constraints:
\[
x^{plan} \models \Phi_{\mathrm{global}}.
\]
The physical execution may still deviate from the plan:
\[
e_{task\_5}^{plan}=49.96s,\qquad e_{task\_5}^{obs}=51.00s.
\]
Therefore,
\[
\Delta e_{task\_5}=e_{task\_5}^{obs}-e_{task\_5}^{plan}=1.04s.
\]
This deviation does not by itself imply a mission violation. The runtime validator must check whether the corresponding runtime properties remain satisfied.

\[
\boxed{\text{Deviation} \neq \text{Violation}}
\]
\end{frame}

\begin{frame}{From Execution Graph to Execution Certificate}
The execution graph represents the planned execution constructed from the SMT solution:
\[
G_E=(V,E).
\]
For runtime validation, we derive an execution certificate:
\[
G_E \longrightarrow C_E,\qquad
C_E=\langle B_E,\Phi_R\rangle.
\]
\begin{itemize}
    \item $B_E$: planned task-robot bindings and timing information;
    \item $\Phi_R$: runtime-checkable properties derived from design-time constraints.
\end{itemize}

\[
\boxed{C_E \text{ defines what must remain valid during execution.}}
\]
\end{frame}

\begin{frame}{From Design-Time Constraints to Runtime Guarantees}
The offline scheduling model is:
\[
\Phi_{\mathrm{global}}
=
\Phi_{\mathrm{assign}}
\land
\Phi_{\mathrm{cap}}
\land
\Phi_{\mathrm{time}}
\land
\Phi_{\mathrm{prec}}
\land
\Phi_{\mathrm{resource}}.
\]
Solving the model produces:
\[
\Phi_{\mathrm{global}}\rightarrow x^{plan}\rightarrow G_E.
\]
Runtime-checkable properties are then derived:
\[
\Phi_R=
\{
\varphi_{\mathrm{prec}},
\varphi_{\mathrm{deadline}},
\varphi_{\mathrm{resource}}
\}.
\]

\[
\boxed{\text{Design-time constraints}\rightarrow\text{runtime guarantees}}
\]
\end{frame}

\begin{frame}{Expected Execution Semantics}
At time $t$, the execution certificate defines the expected execution state:
\[
\boxed{S^{plan}(t)=Expected(C_E,t)}.
\]
For each task $T_i$:
\[
X_i^{plan}(t)
=
\left\langle
status_i^{plan}(t),
r_i^{plan},
s_i^{plan},
e_i^{plan}
\right\rangle.
\]
The planned task status is:
\[
status_i^{plan}(t)=
\begin{cases}
Pending, & t < s_i^{plan},\\
Running, & s_i^{plan}\leq t < e_i^{plan},\\
Completed, & t\geq e_i^{plan}.
\end{cases}
\]
\[
\boxed{Pending\rightarrow Running\rightarrow Completed}
\]
\end{frame}

\begin{frame}{Runtime Events}
During execution, robots generate runtime events:
\[
o_k=
\left\langle
\tau_k,\,
type_k,\,
r_k,\,
T_i,\,
value_k
\right\rangle.
\]
\begin{itemize}
    \item $\tau_k$: timestamp of event $o_k$;
    \item $type_k$: event type, e.g., \texttt{TaskStart}, \texttt{TaskFinish}, or \texttt{RobotFailure};
    \item $r_k$: robot associated with the event;
    \item $T_i$: task associated with the event, if applicable;
    \item $value_k$: observed value carried by the event.
\end{itemize}

For example:
\[
o_k=\langle 51.00s,\texttt{TaskFinish},robot\_4,task\_5,51.00s\rangle.
\]
\end{frame}

\begin{frame}{Observed Execution State}
The runtime event trace is used to reconstruct the observed execution state:
\[
\boxed{S^{obs}(t)=Reconstruct(\mathcal{T}(t))}.
\]
For each task $T_i$:
\[
X_i^{obs}(t)
=
\left\langle
status_i^{obs}(t),\,
r_i^{obs},\,
s_i^{obs},\,
e_i^{obs}
\right\rangle,
\]
where:
\[
status_i^{obs}(t)\in\{Pending,Running,Completed,Failed\}.
\]
The observed state is updated whenever a new runtime event is received:
\[
\boxed{o_k\rightarrow\mathcal{T}(t)\rightarrow S^{obs}(t)}.
\]
\end{frame}

\begin{frame}{Runtime Deviation}
For task $T_i$, the completion-time deviation is:
\[
\Delta e_i=e_i^{obs}-e_i^{plan}.
\]
For the precedence edge extracted from the experimental SMT schedule:
\[
task\_5 \rightarrow task\_7,
\]
a delay of $task\_5$ affects the runtime precedence property:
\[
\Delta e_{task\_5}
\rightarrow
\varphi^{5,7}_{\mathrm{prec}}:
s^{obs}_{task\_7}\geq e^{obs}_{task\_5}.
\]
\[
\boxed{\text{Deviation is interpreted through the affected mission property.}}
\]
\end{frame}

\begin{frame}{Relevance to Mission Constraints}
Only properties that depend on the changed observed state need to be checked:
\[
\Phi_R^{rel}(t)
=
\{\varphi\in\Phi_R\mid \varphi \text{ depends on } S^{obs}(t)\}.
\]
For a completion-time delay in $task\_5$, the relevant precedence property is:
\[
\varphi^{5,7}_{\mathrm{prec}}:
s^{obs}_{task\_7}\geq e^{obs}_{task\_5}.
\]
This avoids treating every timing deviation as a violation.

\[
\boxed{\text{Deviation}\rightarrow\text{relevant runtime property}}
\]
\end{frame}

\begin{frame}{Quantitative Validity Margin}
The validity margin measures the distance from violation.
For a precedence edge $T_i\rightarrow T_j$:
\[
m^{prec}_{i,j}(t)=s^{obs}_{j}-e^{obs}_{i}.
\]
\[
m^{prec}_{i,j}(t)\geq 0
\Rightarrow
\text{Satisfied}
\]
\[
m^{prec}_{i,j}(t)<0
\Rightarrow
\text{Violated}
\]
The sign gives the verdict, and the magnitude gives the severity of the violation.
\end{frame}

\begin{frame}{Runtime Precedence Validation Example}
From the 6r12t SMT schedule, the execution certificate contains:
\[
task\_5 \rightarrow task\_7.
\]
The planned schedule gives:
\[
e^{plan}_{task\_5}=49.96s,\quad
s^{plan}_{task\_7}=53.01s,\quad
m^{plan}_{5,7}=3.05s.
\]
At runtime, the validator checks the observed margin:
\[
m^{obs}_{5,7}=s^{obs}_{task\_7}-e^{obs}_{task\_5}.
\]

\scriptsize
\begin{center}
\begin{tabular}{lccc}
\toprule
Case & $e^{obs}_{task\_5}$ & $s^{obs}_{task\_7}$ & Verdict \\
\midrule
Planned execution & 49.96s & 53.01s & Valid, $m=3.05s$ \\
Delayed but valid & 51.00s & 53.01s & Valid, $m=2.01s$ \\
Runtime violation & 54.50s & 53.01s & Violated, $m=-1.49s$ \\
\bottomrule
\end{tabular}
\end{center}
\end{frame}

\begin{frame}{Proposed Runtime Validation Model}
\[
G_E\rightarrow C_E\rightarrow S^{plan}(t)
\]
\[
\mathcal{T}(t)\rightarrow S^{obs}(t)
\]
\[
\Delta(t)=Diff(S^{plan}(t),S^{obs}(t))
\]
\[
\Phi_R^{rel}(t)=Relevant(\Delta(t),\Phi_R)
\]
\[
\boxed{
(V,D,m)=Validate(S^{obs}(t),\Phi_R^{rel}(t))
}
\]
For the runtime violation above, the validator returns:
\[
V=\text{Violated},\quad
D=\text{precedence}(task\_5,task\_7),\quad
m=-1.49s.
\]
\end{frame}

\begin{frame}{Open Questions for Discussion}
\textbf{Q1. Runtime properties}

Which design-time constraints should be monitored at runtime?

\vspace{0.5cm}
\textbf{Q2. Scope}

Should this paper focus on:
\[
Detection\rightarrow Diagnosis\rightarrow Margin
\]
and leave repair for future work?

\[
\boxed{\text{Deviation}\neq\text{Violation}}
\]
\end{frame}

\end{document}
